\section{Non-Abelian prepared refinement and fusion}\label{res:44-nonabelian-prepared-refinement-and-fusion}
The rooted operator of Section~\ref{sec:tree-refinement-boundary} fixes the
gauge at the root.  For $C_2\le S_3$ it therefore gives the transposition
sector the norm squared $1/3$, and for $C_3\le S_3$ it merges the orbits of
$r$ and $r^2$ (Proposition~\ref{prop:three-refinement-tests}).  This section
keeps a $K$-valued gauge workspace at every vertex, including the root, and
derives a weighted comparison that repairs the first case and states exactly
what the second admits.

The comparisons are prepared gauge boundaries in the sense of
Section~\ref{res:43-prepared-gauge-boundaries-and-composition}.  Such a
boundary admits a subspace $D_e$ of the source physical space, prepares fresh
resources by an owned source run, and supplies an isometry $b_e$ from $D_e$
into the target physical space together with an environment vector $\eta_e$
that does not depend on the input: the raw comparison sends each prepared
input $x$ to $j_Yb_ex\otimes\eta_e$
(Equation~\eqref{eq:43-prepared-gauge-boundaries-and-composition:C8.2}).  The
tag FullJoint requires $\eta_e$ to be invariant under the full target action.
Two boundaries compose on the maximal domain $\{x\in D_f:b_fx\in D_g\}$
(Equation~\eqref{eq:43-prepared-gauge-boundaries-and-composition:C8.12}).  We
write $\mathcal Q_H$ and $\mathcal Q_K$ for the physical spaces of the pure $H$-
and $K$-profiles (written $Q_X$ in
Section~\ref{res:43-prepared-gauge-boundaries-and-composition}).  All optional
fairness, scheduling and independence slots of the profiles and phases in this
chapter are $\AFNoClaim$.

\subsection{The finite gauge-action correspondence}\label{md:44-nonabelian-prepared-refinement-and-fusion:the-finite-gauge-action-correspondence}

Fix a finite connected oriented graph $X$ with $v$ vertices and finite groups
$H\le K$.  Connections are $U\in H^E$, with gauge convention
\begin{equation}
(U^h)_{u\to w}=h_wU_{u\to w}h_u^{-1}.
 \label{eq:44-nonabelian-prepared-refinement-and-fusion:C9.1}
\end{equation}

Let $O_i$ be the full $H^V$ gauge orbits, $n_i=|O_i|$, and $C$ the full $K^V$
orbits meeting $H^E$.  Write $p(i)=C$ when $O_i\subseteq C$, $m_C=|C|$, and
$\kappa_C=|\Stab_{K^V}(U)|$ for $U$ in $C$.  For a connected graph $\kappa_C$
is the cardinality of the simultaneous centralizer of its based chord tuple.
Thus $m_C=|K|^v/\kappa_C$.  We say that $p$ \emph{fuses} two source orbits
$i\ne i'$ when $p(i)=p(i')$.

The comparison carrier contains \textbf{all vertex gauges}, including the root:
\begin{equation}
\mathcal A=H^E\times K^V,\qquad t(U,k)=U^k.
 \label{eq:44-nonabelian-prepared-refinement-and-fusion:C9.2}
\end{equation}

The source profile action is
\begin{equation}
h\cdot(U,k)=(U^h,hkh^{-1}),\quad h\in H^V,
 \qquad t(h\cdot(U,k))=t(U,k)^h.
 \label{eq:44-nonabelian-prepared-refinement-and-fusion:C9.3}
\end{equation}

In particular, a constant source gauge can act nontrivially on the gauge
workspace by conjugation.  Treating that workspace as neutral would lose
the residual symmetry needed in a non-Abelian example.

There is also a full target action on the comparison set $\mathcal A$ by
$(U,k)\mapsto(U,gk)$, with $t(U,gk)=t(U,k)^g$.  Its restriction is not
identical to \eqref{eq:44-nonabelian-prepared-refinement-and-fusion:C9.3}.
The discrepancy is the fibre-preserving transformation
$(U,k)\mapsto(U^h,kh^{-1})$.  Both actions and the fibre relation are retained;
no pullback through a nonexistent inverse of $t$ is asserted.

\subsubsection{Fibre normalization}\label{md:44-nonabelian-prepared-refinement-and-fusion:fibre-normalization}

Put
\begin{equation}
N_C=\sum_{i:p(i)=C}n_i,\qquad \ell_C=N_C\kappa_C.
 \label{eq:44-nonabelian-prepared-refinement-and-fusion:C9.4}
\end{equation}

\begin{proposition}\label{prop:fibre-normalization}
Every $Y\in C$ has exactly $\ell_C$ preimages under $t$.  The weighted
linearization
\begin{equation}
F|U,k\rangle=\ell_{p(i(U))}^{-1/2}|U^k\rangle
 \label{eq:44-nonabelian-prepared-refinement-and-fusion:C9.5}
\end{equation}
is a coisometry onto the raw target saturation $S=K^V\cdot H^E$.  It
intertwines the source action
\eqref{eq:44-nonabelian-prepared-refinement-and-fusion:C9.3} with the embedded
target action.  Its initial projection $F^\dagger F$ averages coefficients
uniformly within each $t$-fibre.
\end{proposition}

\begin{proof}
For a fixed $U$ in $O_i\subseteq C$, the solutions of $U^k=Y$ form a coset of
$\Stab_{K^V}(U)$, of size $\kappa_C$.  Sum over the $N_C$ available source
configurations.  Distinct target rows have disjoint fibres; every row has
$\ell_C$ entries of size $\ell_C^{-1/2}$, so $FF^\dagger=P_S$.  The coefficient
depends only on $C$, and
\eqref{eq:44-nonabelian-prepared-refinement-and-fusion:C9.3} preserves that
class, proving the intertwiner.
\end{proof}

This normalization is not ordinary group averaging.  The unweighted
pushforward of a uniform source and workspace state amplifies fibres.  Nor is
$F$ a deterministic basis pushforward in the sense of
Definition~\ref{def:quantum-gauge-profile-morphism}.  It is an explicitly
weighted correspondence in a prepared boundary of
Section~\ref{res:43-prepared-gauge-boundaries-and-composition}.

\subsection{A legal preparation, with its coefficient boundary exposed}\label{md:44-nonabelian-prepared-refinement-and-fusion:a-legal-preparation-with-its-coefficient-boundary-exposed}

\begin{definition}[Gauge workspace]\label{def:gauge-workspace}
The finite profiles $\AFName{GaugeWorkspace}[H,K,X]$ have the following
fields.  The kinematic field is $H$ connections on canonical $H$ torsors.  At
every vertex $w$, the fresh workspace palette is $K$ with source action
$k_w\mapsto h_wk_wh_w^{-1}$, realized as the standard basis of the
permutation representation $R[K]$; edge coefficient carriers are trivial, with
the usual finite tensor products and evaluation maps.  There is a $0$-cell
coefficient channel at each vertex, supported and anchored there, with this
palette and action and the identity coefficient reduct.  Its
$\AFCoeffIntWit$ has equivariance
\eqref{eq:44-nonabelian-prepared-refinement-and-fusion:C9.3} and a diagonal
loss relation.  Auxiliary state is a singleton; observation and recognition
families are empty; classical Gauss predicates are true.  The group-source
receipts for $H$ and $K$ are declared finite multiplication tables, the
subgroup map is displayed, and the induced-torsor identification is
$[h,a]\mapsto ha$ from $H\times_HK$ to the target canonical $K$ torsor.  No
temporal-origin inheritance is claimed.
\end{definition}

The pure target has $K$ connections; a tagged target is constructed in
Section~\ref{md:44-nonabelian-prepared-refinement-and-fusion:retaining-a-physical-record-changes-the-target-not-the-no-go}.

The full augmented quantum physical carrier is the invariant subspace of
$\mathbb C[H^E\times K^V]$ under
\eqref{eq:44-nonabelian-prepared-refinement-and-fusion:C9.3}, with one
normalized vector per complete orbit.  It is not replaced by the smaller ready
subspace.  The declared fresh resource is $|1\rangle$ at every vertex, fixed
by conjugation.  Thus $j_H|\psi\rangle\otimes|1\rangle^{V}$ is a valid
invariant ready input.

For $d=|K|>1$, let $u=d^{-1/2}\sum_a|a\rangle$ and define the real Householder
matrix $Q_d$ by
\begin{equation}
(Q_d)_{1a}=(Q_d)_{a1}=d^{-1/2},\qquad
 (Q_d)_{ab}=\delta_{ab}-\frac1{d-\sqrt d}\quad(a,b\ne1).
 \label{eq:44-nonabelian-prepared-refinement-and-fusion:C9.6}
\end{equation}

This is $I-2ww^\dagger/\langle w,w\rangle$ for $w=|1\rangle-u$.  Consequently
$Q_d^2=I$ and $Q_d|1\rangle=u$.  Conjugation fixes both $|1\rangle$ and $u$,
so $Q_d$ commutes with it.

The owned phase $\AFName{WorkspacePrepare}[w]$ has, at each input, one
application for every proposed output $a\in K$.  It reads and writes only
$k_w$; its guard is true, its proposal is $k_w:=a$, and its commit law is
finite membership and identity elsewhere.  Legal batches are empty and
singleton; the empty coefficient is zero and the singleton coefficient is
$(Q_d)_{a,k_w}$.  Gauge transport conjugates both labels, preserving the
coefficient.  The trace retains the level, vertex, read, proposal, commit and
coefficient.  The ordered run of the $v$ phases has full raw and physical
unitarity; its reverse is its inverse.

The pure $H$ and $K$ endpoint profiles retain their full raw connection bases,
full normalized gauge-orbit physical bases and inclusions, and an identity run
before any separately named continuation law.  Their other channels are
trivial, with the same declared source receipts and identity reducts.  All
separately owned phases below use the same empty and singleton batch policy,
zero empty coefficient and displayed singleton kernel coefficient.  Every event
retains its phase, read, proposal, commit and coefficient; a preparation
sequence retains each intermediate event.

The run prepares
\begin{equation}
j_H|\psi\rangle\otimes|1\rangle^{V}
   \xrightarrow{\prod_w Q_{d,w}}
 j_H|\psi\rangle\otimes|+_K\rangle^V.
 \label{eq:44-nonabelian-prepared-refinement-and-fusion:C9.7}
\end{equation}

No unknown workspace is reset.  No branch is sampled or copied into an
additional physical record by the trace encoder.

For the concrete groups of orders $2,3,4,6,8,12,24$ used below, the displayed
normalizers and preparation matrices are available in the declared real field
$R=\mathbb Q(\sqrt2,\sqrt3)$, with its positive real embedding into $\mathbb C$
and identity involution.  This is an explicit new scalar choice, not a
conclusion that the dyadic ring $R^{\rm gr}_8$ of
Equation~\eqref{eq:octagonal-ring} already contains the required coefficients.
In particular, \eqref{eq:44-nonabelian-prepared-refinement-and-fusion:C9.6} can
require rational denominators beyond powers of two and three.  General finite
groups require separately sufficient scalars.

\subsection{The normalized physical comparison and its maximal domain}\label{md:44-nonabelian-prepared-refinement-and-fusion:the-normalized-physical-comparison-and-its-maximal-domain}

Let $A$ be the physical map obtained from $F$ and the preparation
\eqref{eq:44-nonabelian-prepared-refinement-and-fusion:C9.7}:
\begin{equation}
F\bigl(j_H|i\rangle\otimes|+_K\rangle^V\bigr)
   =j_K A|i\rangle.
 \label{eq:44-nonabelian-prepared-refinement-and-fusion:C9.8}
\end{equation}

\subsubsection{Weighted fusion}\label{md:44-nonabelian-prepared-refinement-and-fusion:weighted-fusion}

\begin{proposition}\label{prop:weighted-fusion}
For every source orbit label,
\begin{equation}
A|i\rangle=\sqrt{\frac{n_i}{N_{p(i)}}}\,|p(i)\rangle.
 \label{eq:44-nonabelian-prepared-refinement-and-fusion:C9.9}
\end{equation}
Its maximal isometric domain is
\begin{equation}
D=\operatorname{span}\{|s_C\rangle:C\text{ covered}\},\qquad
 |s_C\rangle=\sum_{i:p(i)=C}\sqrt{\frac{n_i}{N_C}}\,|i\rangle.
 \label{eq:44-nonabelian-prepared-refinement-and-fusion:C9.10}
\end{equation}
The vectors $|s_C\rangle$ are orthonormal, $A|s_C\rangle=|C\rangle$, and
$A^\dagger A=P_D$.  Thus the entire source physical space is admitted exactly
when $p$ has no fusion.
\end{proposition}

\begin{proof}
At a fixed $Y\in C$, an input orbit $i$ contributes $n_i\kappa_C$ source pairs,
each with amplitude $1/\sqrt{n_i|K|^v\ell_C}$.  Comparison with the target
amplitude $1/\sqrt{m_C}$ gives
\eqref{eq:44-nonabelian-prepared-refinement-and-fusion:C9.9}.  Each row of $A$
has squared norm one, because its weights sum to $N_C/N_C$.  Different rows
have disjoint source supports.  Therefore $A^\dagger A$ is the projector onto
their adjoint row vectors, which are
\eqref{eq:44-nonabelian-prepared-refinement-and-fusion:C9.10}.  Norm is
preserved exactly on that image.
\end{proof}

The comparison has an exact coherent inverse on this prepared image:
\begin{equation}
F^{\dagger} j_K|C\rangle
     =j_H|s_C\rangle\otimes|+_K\rangle^V.
 \label{eq:44-nonabelian-prepared-refinement-and-fusion:C9.11}
\end{equation}

This makes the role of compression explicit.  The admitted preparation is
already uniform on the fibres, so no projection or postselection is used.  An
arbitrary nonready vector can lose norm under $F$ and is not licensed by the
boundary.  The comparison is still a WITNESS with a legal preceding
preparation, not an automatically installed active quotient SCRIPT.

The source averages are weighted by the numbers of configurations, not by a
uniform count of orbit names.  Write $D_4$ for the dihedral group of order
eight, a Sylow $2$-subgroup of $S_4$.  On one loop, the inclusion
$D_4\subset S_4$ fuses two source classes of sizes one and two.  Their surviving vector has
coefficients $\sqrt{1/3}$ and $\sqrt{2/3}$.  Replacing them by equal weights
fails \eqref{eq:44-nonabelian-prepared-refinement-and-fusion:C9.9}.

\subsection{\texorpdfstring{$C_2\subset S_3$}{C2 in S3}: conjugacy growth without loss of source information}\label{md:44-nonabelian-prepared-refinement-and-fusion:c2s3-conjugacy-growth-without-loss-of-source-information}

Use $H=\{1,(12)\}$ and the square.  Both $H$ orbits have size eight.  The
covered target sectors are the identity and the transposition class
(Table~\ref{tab:c2s3-sectors}).

\begin{table}[htbp]
\centering\small
\renewcommand{\arraystretch}{1.15}
\begin{tabular}{@{}lrrrrr@{}}
\toprule
Sector & Source orbit $n_i$ & Target orbit $m_C$ & Stabilizer $\kappa_C$
  & Fibre $\ell_C$ & Coefficient in $F$\\
\midrule
Identity & $8$ & $216$ & $6$ & $48$ & $1/\sqrt{48}$\\
Transposition & $8$ & $648$ & $2$ & $16$ & $1/4$\\
\bottomrule
\end{tabular}
\caption{The covered sectors of $C_2\subset S_3$ on the square: source and
target orbit sizes, the order $\kappa_C$ of the target stabilizer, the span
fibre size $\ell_C=N_C\kappa_C$ and the coefficient $\ell_C^{-1/2}$ of the
linearization \eqref{eq:44-nonabelian-prepared-refinement-and-fusion:C9.5}.}
\label{tab:c2s3-sectors}
\end{table}

The two target labels are distinct, so $D=\mathcal Q_H$ and
\begin{equation}
A|1\rangle=|[1]_{S_3}\rangle,\qquad
 A|(12)\rangle=|[\text{transposition}]_{S_3}\rangle.
 \label{eq:44-nonabelian-prepared-refinement-and-fusion:C9.12}
\end{equation}

Together with the preparation of
Section~\ref{md:44-nonabelian-prepared-refinement-and-fusion:a-legal-preparation-with-its-coefficient-boundary-exposed}
this supplies a FullJoint prepared boundary
(Section~\ref{res:43-prepared-gauge-boundaries-and-composition}) on the whole original
\textbf{source physical space}.  Its separate environment is the tensor unit.
The augmented source profile still has additional physical states; no
whole-augmented-profile morphism is asserted.

The rooted operator of Section~\ref{sec:tree-refinement-boundary}, built on the
tree charts of Section~\ref{res:40-gauge-fixing-and-refinement-correspondences},
fixes the root gauge and so omits the residual conjugations; it gives the
transposition sector the norm $\sqrt{\rho_a}=1/\sqrt3$.  The root workspace
supplies those conjugations, and the fibre normalization, not a rescaling of
$1/\sqrt3$ to one, restores norm one.  The square calculation enumerates all
$20{,}736=2^4\cdot6^4$ source and workspace pairs and both full target orbits.

\subsection{\texorpdfstring{$C_3\subset S_3$}{C3 in S3}: a fusion boundary and an owned preparation}\label{md:44-nonabelian-prepared-refinement-and-fusion:c3s3-a-genuine-fusion-boundary-and-an-owned-preparation}

Let $r=(123)$.  The three source square orbits have size $27$.  The target
identity orbit has size $216$ and stabilizer six.  The target three-cycle orbit
has size $432$ and stabilizer three; it receives both source labels $r$ and
$r^2$, which $p$ therefore fuses (Figure~\ref{fig:c3s3-fusion}).  In both
target sectors $\ell_C=162$.  Hence
\begin{equation}
A|1\rangle=|[1]\rangle,\quad
 A|r\rangle=2^{-1/2}|[r]\rangle,\quad
 A|r^2\rangle=2^{-1/2}|[r]\rangle.
 \label{eq:44-nonabelian-prepared-refinement-and-fusion:C9.13}
\end{equation}

\begin{figure}[htbp]
\centering
\begin{tikzpicture}[x=1cm,y=1cm,
  vtx/.style={circle,fill=platenavy,inner sep=1.6pt},
  same/.style={draw=platenavy,thick,-{Latex[length=2mm]}},
  moved/.style={draw=platered,line width=1.3pt,-{Latex[length=2mm]}}]
\newcommand{\fusq}[3]{% #1 x shift, #2 style of e_3, #3 label of e_3
  \begin{scope}[shift={(#1,0)}]
    \fill[tealrod!8] (0,0) rectangle (2.3,2.0);
    \draw[same] (0,0)--(2.3,0) node[midway,below,plate label]{$1$};
    \draw[same] (2.3,0)--(2.3,2.0) node[midway,right,plate label]{$1$};
    \draw[same] (2.3,2.0)--(0,2.0) node[midway,above,plate label]{$1$};
    \draw[#2] (0,2.0)--(0,0) node[midway,left,plate label]{#3};
    \foreach \x/\y in {0/0,2.3/0,2.3/2.0,0/2.0}{\node[vtx] at (\x,\y) {};}
    \node[font=\scriptsize] at (1.15,1.0) {$f$};
    \node[font=\scriptsize,anchor=north east] at (0,-.05) {$v_0$};
    \node[font=\scriptsize,anchor=north west] at (2.3,-.05) {$v_1$};
    \node[font=\scriptsize,anchor=south west] at (2.3,2.05) {$v_2$};
    \node[font=\scriptsize,anchor=south east] at (0,2.05) {$v_3$};
  \end{scope}}
\fusq{0}{same}{$r$}
\fusq{6.3}{moved}{\textcolor{platered}{$r^2$}}
\foreach \x/\y in {6.3/0,8.6/0,8.6/2.0,6.3/2.0}{\draw[purplerod,thick] (\x,\y) circle (.18);}
\draw[script arrow] (3.35,1.0)--(5.25,1.0);
\node[font=\scriptsize,align=center,anchor=south] at (4.3,1.1)
  {constant gauge $(12)$\\in $S_3$, not in $C_3$};
\node[font=\scriptsize,text=greenrod!70!black] at (1.15,-.7)
  {$\Hol=r=(123)$};
\node[font=\scriptsize,text=platered] at (7.45,-.7)
  {$\Hol=(12)\,r\,(12)=r^2$};
\node[font=\scriptsize] (sr) at (1.15,-1.55) {$|r\rangle$ in $\mathcal Q_{C_3}$};
\node[font=\scriptsize] (sq) at (7.45,-1.55) {$|r^2\rangle$ in $\mathcal Q_{C_3}$};
\node[font=\scriptsize] (tc) at (4.3,-2.75)
  {$|[r]\rangle$ in $\mathcal Q_{S_3}$, class $\{r,r^2\}$};
\draw[witness arrow] (sr.south east)--(tc.north west)
  node[midway,below left=-1pt,font=\scriptsize]{$A$: $2^{-1/2}$};
\draw[witness arrow] (sq.south west)--(tc.north east)
  node[midway,below right=-1pt,font=\scriptsize]{$A$: $2^{-1/2}$};
\node[font=\scriptsize,align=center,text=plategray] at (4.3,-3.45)
  {admitted domain $D=\operatorname{span}\{|1\rangle,\,(|r\rangle+|r^2\rangle)/\sqrt2\}$};
\end{tikzpicture}
\caption{Fusion of $r=(123)$ and $r^2$ for $C_3\subset S_3$ on the square.
Top: a $C_3$ connection with holonomy $r$ and its image under the constant
gauge change $(12)$, which lies in $S_3$ but not in $C_3$; it conjugates
every link and sends the holonomy to $r^2$.  Source: $C_3$ connections, whose
holonomy is gauge invariant, so $r$ and $r^2$ label different source orbits.
Target: $S_3$ connections, where the two orbits lie in one class.  Law: the
physical comparison $A$ of Equation~\eqref{eq:44-nonabelian-prepared-refinement-and-fusion:C9.8},
on the same square; bottom: $A|r\rangle=A|r^2\rangle=2^{-1/2}|[r]\rangle$ by
Equation~\eqref{eq:44-nonabelian-prepared-refinement-and-fusion:C9.13}, so $A$
is isometric only on the domain $D$ of
Equation~\eqref{eq:44-nonabelian-prepared-refinement-and-fusion:C9.14}.}
\label{fig:c3s3-fusion}
\end{figure}

The prescribed map sending each source basis state to its normalized target
class cannot be an isometry: it sends the orthogonal states $|r\rangle$ and
$|r^2\rangle$ to the same state.  The canonical comparison instead has exact
initial domain
\begin{equation}
D=\operatorname{span}\{|1\rangle,(|r\rangle+|r^2\rangle)/\sqrt2\}.
 \label{eq:44-nonabelian-prepared-refinement-and-fusion:C9.14}
\end{equation}

Its antisymmetric complement maps to zero.  This is a restriction of the
admitted boundary, not a claim that an antisymmetric source state ceases to be
physical.

A source preparation supplies
\eqref{eq:44-nonabelian-prepared-refinement-and-fusion:C9.14}.  The owned phase
$\AFName{ClassMix}$ fixes flux $1$ and applies the ordinary $2\times2$ Hadamard
to the flux labels $r,r^2$.  In raw square coordinates it reads all links,
preserves the three tree links, and proposes the chord $30$ that gives the
chosen nonidentity output flux.  If $T=U_{23}U_{12}U_{01}$ and the proposed
holonomy is $a$, set $U'_{30}=aT^{-1}$.  The identity-flux branch is the
identity application.  The other two branches have Hadamard coefficients
$\pm1/\sqrt2$.

Guards specify these alternatives; write support is the chord, affected
support is the square, and commit and trace data retain the displayed
assignment and coefficient.  $C_3$ is Abelian, so the flux is invariant and the
proposal commutes with every source gauge.  The matrix is the identity on
identity flux plus a Hadamard on each fixed tree-coordinate block.  With the
gauge workspace present it acts as the identity on that workspace.

Starting in the ready code $\operatorname{span}\{|1\rangle,|r\rangle\}$, first
run $\AFName{ClassMix}$, then the four workspace preparations of
Section~\ref{md:44-nonabelian-prepared-refinement-and-fusion:a-legal-preparation-with-its-coefficient-boundary-exposed}.
Equation~\eqref{eq:44-nonabelian-prepared-refinement-and-fusion:C9.8} maps this
ready code isometrically onto the two covered target classes.  The other ready
input $|r^2\rangle$ is sent by $\AFName{ClassMix}$ to the rejected antisymmetric
sector.  This is a second, nontrivial prepared boundary; its admitted input is
smaller than the three-dimensional source physical space.

\subsection{Retaining a physical record changes the target, not the no-go}\label{md:44-nonabelian-prepared-refinement-and-fusion:retaining-a-physical-record-changes-the-target-not-the-no-go}

To preserve all source labels, give the target a neutral physical register
$i$ identifying the source $H$-orbit.

\begin{definition}[Tagged target]\label{def:tagged-target}
The \emph{tagged target} is the pure $K$-profile with one more finite neutral
coefficient channel at the root, supported and anchored there, with standard
basis $R[I]$ indexed by the set $I$ of source orbit labels, trivial gauge
action and identity reduct.  Auxiliary state remains a singleton.  The target
retains the full physical basis $|C,i\rangle$, including pairs outside the
comparison image, and an identity run.  Its span endpoint and linearization
are
\begin{equation}
t_{\rm tag}(U,k)=(U^k,i(U)),\qquad
 \ell_i=n_i\kappa_{p(i)},\qquad
 F_{\rm tag}|U,k\rangle=\ell_i^{-1/2}|U^k,i(U)\rangle.
 \label{eq:44-nonabelian-prepared-refinement-and-fusion:C9.15}
\end{equation}
\end{definition}

The fibres of $t_{\rm tag}$ no longer merge distinct source orbit labels.  The
counting in the proof of Proposition~\ref{prop:weighted-fusion} gives
\begin{equation}
F_{\rm tag}(j_H|i\rangle\otimes|+_K\rangle^V)
     =j_K|p(i)\rangle\otimes|i\rangle.
 \label{eq:44-nonabelian-prepared-refinement-and-fusion:C9.16}
\end{equation}

This is an isometry on the whole source physical space, including fusion
examples, \textbf{into the enlarged physical target}.  The labels of the
register persist.  The register is part of the target profile of the boundary
map $b_e$, not a fixed independent environment $\eta_e$ in the system-only
factorization of Section~\ref{res:43-prepared-gauge-boundaries-and-composition}.

For a superposition of $r$ and $r^2$, the joint record retains their relative
phase.  Tracing that record gives the same system state for the symmetric and
antisymmetric superpositions.  Thus
\eqref{eq:44-nonabelian-prepared-refinement-and-fusion:C9.16} does not repair
an isometry to the original system-only target.  Nor does it assert an actual
outcome or an observer interpretation for that register.

The normalizer in \eqref{eq:44-nonabelian-prepared-refinement-and-fusion:C9.15}
differs from \eqref{eq:44-nonabelian-prepared-refinement-and-fusion:C9.5} when
fusion occurs.  One cannot attach an external record and keep using the
untagged comparison $F$ without rechecking its domain.  $F$ preserves norm on
vectors constant within its complete fibres; physical copies distinguishing
members of a fibre can destroy that condition.  For example, copy every pair
$(U,k)$ into an orthogonal external register before applying the original $F$
to the source.  On either normalized $C_2$ square sector the resulting squared
norm is $1/\ell_C$, namely $1/48$ or $1/16$; that recorded input is outside the
admitted domain.  Proof and history witnesses are retained separately and do
not create such copies.  No record is erased through a forgetful map.

\subsection{Continuation laws and the locality boundary}\label{md:44-nonabelian-prepared-refinement-and-fusion:continuation-laws-and-the-locality-boundary}

Also define $\AFName{ClassSign}[\chi]$, for a specified $\pm1$ function $\chi$
on target conjugacy-orbit labels.  The source coefficient is $\chi(p(i(U)))$
and the target coefficient is $\chi(C(Y))$.  The application is the identity,
with read support all connections, no writes, true guard, and a singleton
weight equal to that sign; empty batches have zero weight.  Its event retains
the read and sign.  Gauge covariance, closed unitarity and preservation of the
ready fusion code follow immediately from constancy on each covered target
class.  It commutes with workspace preparation.

Coefficientwise, $F\,\AFName{ClassSign}_H=\AFName{ClassSign}_K\,F$.  The
history comparison sends the source sign event at $U$ to the target sign event
at $U^k$, retaining the span witness $(U,k)$ and the normalizing coefficient.
It preserves the complete source and target provenance; it does not identify
different source events merely because their target sign agrees.  This is the
selected law-qualified family for the untagged boundary.  A source sign that
distinguishes fused classes does not descend to this family; the tagged target
can instead read its retained register $i$.

A naive noncentral target link write is not gauge covariant.  For example,
left multiplication of one $S_3$ link by $(12)$ fails covariance under a
constant three-cycle gauge.  The cyclic native write rule cannot simply be
reused with the same name.

There is a separately constructible finite coherent alternative.  For
normalized target orbit vectors $u_0,u_t$ of identity and transposition,
\begin{equation}
K_{\rm swap}=I-|u_0\rangle\langle u_0|-|u_t\rangle\langle u_t|
       +|u_t\rangle\langle u_0|+|u_0\rangle\langle u_t|
 \label{eq:44-nonabelian-prepared-refinement-and-fusion:C9.17}
\end{equation}
is a raw unitary commuting with every gauge transformation.  The owned phase
$\AFName{OrbitSwap}$ has one application for each output configuration in the
supported row of this kernel, reads and writes the full graph, and uses the
displayed matrix coefficient.  Its guards, proposals and commits are finite
membership; its traces keep all coefficients, including negative same-orbit
terms.  On other orbits it is the identity.  Its complete physical kernel swaps
the two selected labels and fixes the others.

This is a full-support finite law, not a single-link write.  It satisfies the
corresponding physical kernel square for
\eqref{eq:44-nonabelian-prepared-refinement-and-fusion:C9.12}, but that
equality does not supply a conservative comparison of the complete histories
of the cyclic write.  Such a stronger law interface is not claimed here.
Direct computation checks the one-loop raw matrix; the projector formula
\eqref{eq:44-nonabelian-prepared-refinement-and-fusion:C9.17} proves the
statement on the square.

\subsection{Composition: one positive tower and one readiness failure}\label{md:44-nonabelian-prepared-refinement-and-fusion:composition-one-positive-tower-and-one-real-readiness-failure}

For $C_2\subset S_3\subset S_3\times C_2$, use the subgroup in the first factor
at each stage.  Both orbit maps have no fusion.  The prepared physical maps
therefore compose as the direct inclusion of labels.  Each stage keeps its
independently supplied gauge workspace, preparation events, span and
normalizer in the staged comparison.  The second stage accepts the three-cycle
class absent from the first landing.  The finite checks use one loop; the
no-fusion statement for the first inclusion also holds for simultaneous
tuples, since conjugating a nonidentity tuple entry $(12)$ back into $C_2$
fixes it.  The direct-product second inclusion creates no new conjugacy
identifications.

In contrast, take $C_2\subset V_4\subset S_4$, where $V_4$ is the Klein
subgroup of double transpositions.  The first map sends its nonidentity basis
vector to one of the three nonidentity $V_4$ labels.  The unadapted canonical
second map accepts only their uniform superposition and sends that one label
with amplitude $1/\sqrt3$.  The direct $C_2\subset S_4$ comparison has no
fusion and sends it with amplitude one.  Thus
\begin{equation}
A_{V_4,S_4}A_{C_2,V_4}\ne A_{C_2,S_4}.
 \label{eq:44-nonabelian-prepared-refinement-and-fusion:C9.18}
\end{equation}

The exact staged ready domain, without an adapter, is only the identity line.
The apparent norm squared $1/3$ on the other line is failed readiness, not a
success probability for a licensed refinement run.

An explicit source preparation repairs that selected path.  On the three
nonidentity $V_4$ flux labels, use $Q_3$ from
\eqref{eq:44-nonabelian-prepared-refinement-and-fusion:C9.6}, fixing identity
flux.  This phase $\AFName{V4ClassMix}$ has the same finite chord-proposal
construction as $\AFName{ClassMix}$, with three alternatives and the $Q_3$
coefficients.  $V_4$ is Abelian, so it is an owned covariant unitary.  It maps
the particular incoming double-transposition label to the required uniform
combination.  With this extra preparation, the two-dimensional incoming code
is admitted and the physical staged map equals the direct map.  The added
process and its events remain part of the path; they are not inferred from the
subgroup inclusions alone.

This is a group-theoretic instance of the guarded composition of prepared
boundaries (Section~\ref{res:43-prepared-gauge-boundaries-and-composition},
Equation~\eqref{eq:43-prepared-gauge-boundaries-and-composition:C8.12}).  It
also explains why a universal refinement tower cannot be inferred from
successful individual normalized comparisons.
Table~\ref{tab:inclusion-summary} compares these towers with the single
inclusions of this chapter.
